\section{问题一：给定方案下的系统状态}

\subsection{问题目标与建模思路}

在给定 II 型锚链、锚链总长和重物球质量的条件下，分别考察风速为
\(12\,\mathrm{m/s}\) 与 \(24\,\mathrm{m/s}\) 时的系统状态。根据第
\ref{sec:assumptions}节，四节钢管均为完全密闭的刚性圆柱形构件，重物球通过
质量可忽略的竖直短绳与钢桶相连。以浮标吃水深度 \(\delta\) 为自变量，按照浮标、
钢管、钢桶--重物球及锚链的受力平衡关系逐段求解，并由各构件的竖直投影和水平
投影分别计算海水深度与浮标游动半径。

图\ref{fig:q1-freebody}给出三个基本受力单元。图中的 \(T_i\) 表示相邻构件
之间的张力，\(F_i\) 与 \(G_i\) 分别表示构件所受浮力和重力；下标
\(\mathrm d\) 与 \(\mathrm s\) 分别表示钢桶和重物球。钢桶--重物球单元中，
重物球的作用通过 \(G_{\mathrm s}-F_{\mathrm s}\) 传递至钢桶下端。

\begin{figure}[H]
  \centering
  \begin{subfigure}[b]{0.30\textwidth}
    \centering
    \includesvg[width=\linewidth]{figures/q1_buoy_force}
    \caption{浮标受力}
  \end{subfigure}\hfill
  \begin{subfigure}[b]{0.30\textwidth}
    \centering
    \includesvg[width=\linewidth]{figures/q1_pipe_force}
    \caption{钢管通用受力}
  \end{subfigure}\hfill
  \begin{subfigure}[b]{0.30\textwidth}
    \centering
    \includesvg[width=\linewidth]{figures/q1_bucket_ball_force}
    \caption{钢桶--重物球单元}
  \end{subfigure}
  \caption{问题一关键受力示意图}
  \label{fig:q1-freebody}
\end{figure}

\subsection{浮标与钢管的受力平衡}

浮标为半径 \(r\)、直径 \(D\) 的圆柱体。吃水深度为 \(\delta\) 时，其浮力和
风力分别为
\begin{equation}
  F_{\mathrm b}=\rho_{\text{海水}}g\pi r^2\delta,
  \label{eq:q1-buoyancy}
\end{equation}
\begin{equation}
  F_{\mathrm w}=0.625\times(2-\delta)D\times v^2.
  \label{eq:q1-wind-load}
\end{equation}
由图\ref{fig:q1-freebody}中浮标的水平和竖直受力平衡，得到
\begin{equation}
  F_{\mathrm w}=T_1\cos\theta_1,
  \label{eq:q1-buoy-horizontal}
\end{equation}
\begin{equation}
  G_{\mathrm b}+T_1\sin\theta_1=F_{\mathrm b}.
  \label{eq:q1-buoy-vertical}
\end{equation}

以第 \(i\) 节钢管为对象，相邻两端的力平衡写为
\begin{equation}
  \begin{cases}
    T_{i+1}\sin\theta_{i+1}=T_i\sin\theta_i,\\[4pt]
    T_{i+1}\cos\theta_{i+1}
    =T_i\cos\theta_i+F_i-G_i,
  \end{cases}
  \qquad i=1,2,3,4.
  \label{eq:q1-pipe-force-balance}
\end{equation}
对钢管两端结点分别取矩，有
\begin{equation}
  (G_i-F_i)\sin\varphi_i\times\frac{1}{2}L_i
  =T_{i+1}\sin(\theta_{i+1}-\varphi_i)\times L_i,
  \label{eq:q1-pipe-moment-upper}
\end{equation}
\begin{equation}
  (G_i-F_i)\sin\varphi_i\times\frac{1}{2}L_i
  =T_i\sin(\varphi_i-\theta_i)\times L_i.
  \label{eq:q1-pipe-moment-lower}
\end{equation}
由此得到钢管姿态角的递推表达式
\begin{equation}
  \tan\varphi_i=
  \frac{2T_i\cos\theta_i-(G_i-F_i)}
  {2T_i\sin\theta_i}.
  \label{eq:q1-pipe-angle}
\end{equation}

\subsection{钢桶、重物球与锚链段的受力平衡}

钢桶与重物球单元的两端力平衡为
\begin{equation}
  T_6\sin\theta_6=T_5\sin\theta_5,
  \label{eq:q1-bucket-horizontal}
\end{equation}
\begin{equation}
  T_6\sin\theta_6
  =T_5\sin\theta_5-(G_{\mathrm s}-F_{\mathrm s})-(G_{\mathrm d}-F_{\mathrm d}).
  \label{eq:q1-bucket-vertical}
\end{equation}
分别对钢桶两端取矩，得到
\begin{equation}
  \frac{1}{2}L_5(G_{\mathrm d}-F_{\mathrm d})\sin\varphi_5
  =T_5\sin(\varphi_5-\theta_5)L_5,
  \label{eq:q1-bucket-moment-upper}
\end{equation}
\begin{equation}
  \frac{1}{2}L_5(G_{\mathrm d}-F_{\mathrm d})\sin\varphi_5
  -(G_{\mathrm s}-F_{\mathrm s})L_5\sin\varphi_5
  =L_5T_6\sin(\theta_6-\varphi_5).
  \label{eq:q1-bucket-moment-lower}
\end{equation}

锚链段沿受力链路继续采用相同的逐节平衡关系，编号为
\(i=6,7,8,\ldots,n\)。将各节构件的竖直投影相加，海水深度为
\begin{equation}
  H=\sum_{i=1}^{n}L_i\cos\theta_i+h.
  \label{eq:q1-water-depth}
\end{equation}
相应地，浮标以锚所在位置为圆心的游动半径为
\begin{equation}
  R=\sum_{i=1}^{n}L_i\sin\theta_i.
  \label{eq:q1-swing-radius}
\end{equation}

\subsection{数值求解、脚本与可复现输出}

上述模型由可执行脚本 \texttt{src/q1/q1\_static\_equilibrium.py} 实现。该脚本不改写原始
数据：先在 \([0.670,0.850]\,\mathrm m\) 内取 \(\delta\)，计算
\(F_{\mathrm w}\)、四节钢管与钢桶姿态以及锚链状态，再由
式\eqref{eq:q1-water-depth}计算残差 \(f(\delta)=H(\delta;v)-18\)。对每个风速
执行 80 次二分迭代，取得 \(f(\delta)=0\) 的数值解；随后按
式\eqref{eq:q1-swing-radius}计算游动半径。仓库根目录下的运行命令为
\verb|python src/q1/q1_static_equilibrium.py|。

每次运行均自动写出：吃水--水深曲线及两种风速的锚链坐标到
\texttt{data/processed/}，最终结果 CSV/\LaTeX{} 表到
\texttt{outputs/q1/tables/}，图件到 \texttt{outputs/q1/figures/}，并将边界、
水深残差和锚链长度残差写入 \texttt{outputs/q1/logs/q1\_static\_checks.txt}。
因此，以下图表均由脚本输出读取生成，而非手工填入。

\subsection{求解结果与核验}

图\ref{fig:q1-depth-draft}显示 \(H(\delta;v)\) 与目标水深的交点；图
\ref{fig:q1-chain-shapes}展示同一求解结果下的锚链形状。两图均优先由
\texttt{paper/figures/} 中的 TikZ 源引入，数据路径保持相对路径。

\begin{figure}[H]
  \centering
  \includestandalone[mode=tex,width=0.92\linewidth]{figures/q1_depth_vs_draft}
  \caption{不同风速下浮标吃水与模型水深的关系}
  \label{fig:q1-depth-draft}
\end{figure}

\begin{figure}[H]
  \centering
  \includestandalone[mode=tex,width=0.82\linewidth]{figures/q1_chain_shapes}
  \caption{不同风速下锚链的静力平衡形状}
  \label{fig:q1-chain-shapes}
\end{figure}

\begin{table}[H]
  \centering
  \caption{问题一两种风速下的求解结果}
  \label{tab:q1-static-results}
  \small
  % 由 src/q1/q1_static_model.m 自动生成；请勿手工修改。
\begin{tabular}{lrr}
\toprule
指标 & 12 m/s & 24 m/s \\ 
\midrule
第1节钢管与海平面竖直方向夹角（度） & 1.160 & 4.413 \\ 
第2节钢管与海平面竖直方向夹角（度） & 1.168 & 4.441 \\ 
第3节钢管与海平面竖直方向夹角（度） & 1.175 & 4.470 \\ 
第4节钢管与海平面竖直方向夹角（度） & 1.184 & 4.499 \\ 
钢桶与海平面竖直方向夹角（度） & 1.202 & 4.566 \\ 
最后一条锚链与海平面水平方向夹角（度） & 0.000 & 4.471 \\ 
浮标的吃水深度（m） & 0.683 & 0.697 \\ 
浮标的游动半径（m）（以锚所在位置为圆心） & 14.654 & 17.780 \\ 
\bottomrule
\end{tabular}

\end{table}

在 \(12\,\mathrm{m/s}\) 下，锚链有 \(6.249\,\mathrm m\) 卧底；在
\(24\,\mathrm{m/s}\) 下，锚链全长悬空。核验日志给出的两种工况水深残差分别为
\(7.105\times10^{-15}\,\mathrm m\) 与 \(3.553\times10^{-14}\,\mathrm m\)，锚链
长度残差均为零，且两解的吃水均在浮标高度范围内。这说明在本节既定的准静态与
构件密闭假设下，数值闭合满足设定水深和链长约束。
